```latex
\documentclass[12pt]{article}
\usepackage[utf8]{inputenc}
\usepackage{amsmath, amssymb, graphicx, hyperref}
\usepackage{geometry}
\geometry{a4paper, margin=1in}
\title{SHOA-Praktik: GRAPE Framework -- A Universal Resilience Architecture for Computing Systems}
\author{Konstantin Fedotov}
\date{\today}

\begin{document}

\maketitle

\begin{abstract}
We present the SHOA-Praktik GRAPE Framework, the second practical module of the SHOA research program. Building on the SHOA-Neuro, SHOA-GRAPE, and NeuroSHOA preprints, the framework provides a universal resilience architecture for any computing system. It implements the Spiked Voting protocol across 7 independent computational channels, a Gated Redundancy mechanism with ``pre-failure memory'', and a software Grape pulse for automated recovery. The framework includes a real-time monitoring dashboard, an integration API for major AI and quantum computing platforms, and a stress-test simulator. We provide four formal validation criteria: recovery time, state selection accuracy, fault resilience, and cross-platform compatibility. The GRAPE Framework transforms the SHOA principles of collective voting and gradient-optimized recovery into a deployable software layer for critical infrastructure.
\end{abstract}

\section{Introduction}
Modern computing systems, from data centers to autonomous drones, face an escalating threat landscape: hardware failures, software errors, and cyberattacks. Classical fault-tolerance mechanisms rely on static redundancy and manual intervention, which are too slow and inflexible for real-time recovery.

The SHOA project has established a superior alternative: dynamic, collective decision-making via the Spiked Voting protocol and adaptive recovery through a Grape pulse. The GRAPE Framework is the software embodiment of these principles. It is designed to be integrated into any existing system with minimal overhead, providing a universal resilience layer that learns from every fault.

\section{Design and Architecture}
\subsection{System Components}
The GRAPE Framework consists of five integrated subsystems:
\begin{itemize}
    \item \textbf{Spiked Voting Library:} A collection of configurable voting algorithms (majority, weighted, dynamic) that make collective decisions across up to 7 independent channels in under 1 millisecond.
    \item \textbf{Gated Redundancy Module:} Manages 7 independent computational channels, each periodically saving its ``pre-failure memory'' (a complete snapshot of its stable state). Faulty channels are automatically isolated, and reserve nodes are activated.
    \item \textbf{State Monitoring System:} Tracks CPU/RAM load, response delays, computation errors, and data integrity in real time with a configurable polling frequency of 10--100 Hz.
    \item \textbf{Integration API:} Provides connectivity with TensorFlow, PyTorch, Qiskit, and Kubernetes via REST, gRPC, and MQTT protocols.
    \item \textbf{Management Dashboard:} A web-based interface for visualizing channel status, the resilience graph, failure/recovery history, and for sending alerts via email, Telegram, or SMS.
\end{itemize}

\subsection{Operational Workflow}
The recovery cycle follows a strict sequence:
\begin{enumerate}
    \item \textbf{Fault Detection:} The monitoring system registers a deviation (error, delay, lost connection) exceeding the configured threshold.
    \item \textbf{Voting:} The Spiked Voting Library is activated. All healthy channels vote for the most stable state based on their pre-failure memory.
    \item \textbf{Grape Pulse:} A software-defined Grape pulse triggers recovery. The faulty channel is isolated, a reserve node is connected, and the system restarts from the voted state.
    \item \textbf{Verification:} The monitoring system confirms stability, and the system returns to normal mode. Data on the fault and recovery is stored for future learning.
\end{enumerate}

\section{Formal Validation Criteria}
To ensure that the GRAPE Framework provides genuine SHOA-level resilience, we impose the following mandatory criteria.

\subsection{Criterion I: Recovery Time}
\begin{itemize}
    \item \textbf{Definition:} The system must recover from a software fault in under 10 ms and from a simulated hardware fault in under 1 second.
    \item \textbf{Measurement:} Inject a fault (e.g., kill a process or corrupt a memory block) and measure the time from detection to full restoration of service.
    \item \textbf{Target:} $t_{\text{soft}} < 10$ ms; $t_{\text{hard}} < 1$ s.
\end{itemize}

\subsection{Criterion II: State Selection Accuracy}
\begin{itemize}
    \item \textbf{Definition:} The Spiked Voting must select the correct (most stable) state from the pre-failure memory with an accuracy exceeding 98\%.
    \item \textbf{Measurement:} Run 1000 fault-recovery cycles with a known correct state and count the number of incorrect selections.
    \item \textbf{Target:} Accuracy $> 98\%$.
\end{itemize}

\subsection{Criterion III: Fault Resilience Under Load}
\begin{itemize}
    \item \textbf{Definition:} The system must maintain 99.99\% uptime under a continuous stream of simulated attacks and random faults over a 72-hour stress test.
    \item \textbf{Measurement:} Deploy the framework on a 7-node cluster and run a script that injects faults at random intervals. Measure total downtime.
    \item \textbf{Target:} Uptime $> 99.99\%$ (less than 8.6 seconds of downtime over 72 hours).
\end{itemize}

\subsection{Criterion IV: Cross-Platform Compatibility}
\begin{itemize}
    \item \textbf{Definition:} The framework must successfully integrate with at least three major AI or quantum computing platforms.
    \item \textbf{Measurement:} Deploy a simple model (e.g., a TensorFlow classifier) and a Qiskit circuit, and verify that the framework correctly monitors, detects faults, and triggers recovery for both.
    \item \textbf{Target:} Successful integration with TensorFlow, PyTorch, and Qiskit.
\end{itemize}

\section{Conclusion}
The SHOA-Praktik GRAPE Framework is the first universal, software-defined resilience architecture based on the principles of quantum emergence and collective voting. By satisfying the four formal criteria outlined here, any organization can independently verify that their critical systems are protected by a SHOA-compliant resilience layer.

\end{document}
